% Traducción española de la tabla de materias de EGA IV, cuarta parte, pp. 341--343.
\clearpage
\oldpage[IV]{341}
\phantomsection
\addcontentsline{toc}{section}{Índice}
\section*{ÍNDICE}
\markboth{ÍNDICE}{ÍNDICE}

\begingroup
\small
\makeatletter
\newcommand{\EGAIVcontentsline}[3]{%
  \@dottedtocline{1}{0em}{3.5em}{\textbf{#1}\quad #2}{#3}}
\makeatother

\begin{center}\textbf{CAPÍTULO IV. --- Estudio local de los esquemas y de los morfismos de esquemas (fin)}\end{center}

\EGAIVcontentsline{\S16}{Invariantes diferenciales. Morfismos diferencialmente lisos}{5}
\EGAIVcontentsline{{\texorpdfstring{\protect\EGAReadLink{IV.16.1}{16.1}}{16.1}}}{Invariantes normales de una inmersión}{5}
\EGAIVcontentsline{{\texorpdfstring{\protect\EGAReadLink{IV.16.2}{16.2}}{16.2}}}{Propiedades functoriales de los invariantes normales de una inmersión}{9}
\EGAIVcontentsline{{\texorpdfstring{\protect\EGAReadLink{IV.16.3}{16.3}}{16.3}}}{Invariantes diferenciales fundamentales de un morfismo de preesquemas}{14}
\EGAIVcontentsline{{\texorpdfstring{\protect\EGAReadLink{IV.16.4}{16.4}}{16.4}}}{Propiedades functoriales de los invariantes diferenciales}{16}
\EGAIVcontentsline{{\texorpdfstring{\protect\EGAReadLink{IV.16.5}{16.5}}{16.5}}}{Haces y fibrados tangentes relativos; derivaciones}{27}
\EGAIVcontentsline{{\texorpdfstring{\protect\EGAReadLink{IV.16.6}{16.6}}{16.6}}}{Haces de $p$-diferenciales y diferencial exterior}{34}
\EGAIVcontentsline{{\texorpdfstring{\protect\EGAReadLink{IV.16.7}{16.7}}{16.7}}}{Los $\sh{P}^{n}_{X/S}(\sh{F})$}{36}
\EGAIVcontentsline{{\texorpdfstring{\protect\EGAReadLink{IV.16.8}{16.8}}{16.8}}}{Operadores diferenciales}{39}
\EGAIVcontentsline{{\texorpdfstring{\protect\EGAReadLink{IV.16.9}{16.9}}{16.9}}}{Inmersiones regulares y cuasirregulares}{46}
\EGAIVcontentsline{{\texorpdfstring{\protect\EGAReadLink{IV.16.10}{16.10}}{16.10}}}{Morfismos diferencialmente lisos}{51}
\EGAIVcontentsline{{\texorpdfstring{\protect\EGAReadLink{IV.16.11}{16.11}}{16.11}}}{Operadores diferenciales sobre un $S$-preesquema diferencialmente liso}{53}
\EGAIVcontentsline{{\texorpdfstring{\protect\EGAReadLink{IV.16.12}{16.12}}{16.12}}}{Caso de característica cero: criterio jacobiano para los morfismos diferencialmente lisos}{55}

\EGAIVcontentsline{\S17}{Morfismos lisos, morfismos no ramificados, morfismos étale}{56}
\EGAIVcontentsline{{\texorpdfstring{\protect\EGAReadLink{IV.17.1}{17.1}}{17.1}}}{Morfismos formalmente lisos, morfismos formalmente no ramificados, morfismos formalmente étale}{56}
\EGAIVcontentsline{{\texorpdfstring{\protect\EGAReadLink{IV.17.2}{17.2}}{17.2}}}{Propiedades generales de los diferenciales}{59}
\EGAIVcontentsline{{\texorpdfstring{\protect\EGAReadLink{IV.17.3}{17.3}}{17.3}}}{Morfismos lisos, morfismos no ramificados, morfismos étale}{61}
\EGAIVcontentsline{{\texorpdfstring{\protect\EGAReadLink{IV.17.4}{17.4}}{17.4}}}{Caracterizaciones de los morfismos no ramificados}{65}
\EGAIVcontentsline{{\texorpdfstring{\protect\EGAReadLink{IV.17.5}{17.5}}{17.5}}}{Caracterizaciones de los morfismos lisos}{67}
\EGAIVcontentsline{{\texorpdfstring{\protect\EGAReadLink{IV.17.6}{17.6}}{17.6}}}{Caracterizaciones de los morfismos étale}{70}
\EGAIVcontentsline{{\texorpdfstring{\protect\EGAReadLink{IV.17.7}{17.7}}{17.7}}}{Propiedades de descenso y de paso al límite}{72}
\EGAIVcontentsline{{\texorpdfstring{\protect\EGAReadLink{IV.17.8}{17.8}}{17.8}}}{Criterios de lisitud y de no ramificación por fibras}{79}
\EGAIVcontentsline{{\texorpdfstring{\protect\EGAReadLink{IV.17.9}{17.9}}{17.9}}}{Morfismos étale e inmersiones abiertas}{79}
\EGAIVcontentsline{{\texorpdfstring{\protect\EGAReadLink{IV.17.10}{17.10}}{17.10}}}{Dimensión relativa de un preesquema liso sobre otro}{81}
\EGAIVcontentsline{{\texorpdfstring{\protect\EGAReadLink{IV.17.11}{17.11}}{17.11}}}{Morfismos lisos de preesquemas lisos}{82}
\EGAIVcontentsline{{\texorpdfstring{\protect\EGAReadLink{IV.17.12}{17.12}}{17.12}}}{Subpreesquemas lisos de un preesquema liso. Morfismos lisos y morfismos diferencialmente lisos}{85}

\oldpage[IV]{342}
\EGAIVcontentsline{{\texorpdfstring{\protect\EGAReadLink{IV.17.13}{17.13}}{17.13}}}{Morfismos transversales}{89}
\EGAIVcontentsline{{\texorpdfstring{\protect\EGAReadLink{IV.17.14}{17.14}}{17.14}}}{Caracterizaciones locales e infinitesimales de los morfismos lisos, de los morfismos no ramificados y de los morfismos étale}{98}
\EGAIVcontentsline{{\texorpdfstring{\protect\EGAReadLink{IV.17.15}{17.15}}{17.15}}}{Caso de los preesquemas sobre un cuerpo de base}{99}
\EGAIVcontentsline{{\texorpdfstring{\protect\EGAReadLink{IV.17.16}{17.16}}{17.16}}}{Cuasi-secciones de morfismos planos o lisos}{105}

\EGAIVcontentsline{\S18}{Complementos sobre los morfismos étale. Anillos locales henselianos y anillos estrictamente locales}{109}
\EGAIVcontentsline{{\texorpdfstring{\protect\EGAReadLink{IV.18.1}{18.1}}{18.1}}}{Una equivalencia notable de categorías}{109}
\EGAIVcontentsline{{\texorpdfstring{\protect\EGAReadLink{IV.18.2}{18.2}}{18.2}}}{Recubrimientos étale}{111}
\EGAIVcontentsline{{\texorpdfstring{\protect\EGAReadLink{IV.18.3}{18.3}}{18.3}}}{Álgebras finitas y étale}{114}
\EGAIVcontentsline{{\texorpdfstring{\protect\EGAReadLink{IV.18.4}{18.4}}{18.4}}}{Estructura local de los morfismos no ramificados y de los morfismos étale}{118}
\EGAIVcontentsline{{\texorpdfstring{\protect\EGAReadLink{IV.18.5}{18.5}}{18.5}}}{Anillos locales henselianos}{125}
\EGAIVcontentsline{{\texorpdfstring{\protect\EGAReadLink{IV.18.6}{18.6}}{18.6}}}{Henselización}{135}
\EGAIVcontentsline{{\texorpdfstring{\protect\EGAReadLink{IV.18.7}{18.7}}{18.7}}}{Henselización y anillos excelentes}{142}
\EGAIVcontentsline{{\texorpdfstring{\protect\EGAReadLink{IV.18.8}{18.8}}{18.8}}}{Anillos estrictamente locales y henselización estricta}{144}
\EGAIVcontentsline{{\texorpdfstring{\protect\EGAReadLink{IV.18.9}{18.9}}{18.9}}}{Fibras formales de los anillos noetherianos henselianos}{150}
\EGAIVcontentsline{{\texorpdfstring{\protect\EGAReadLink{IV.18.10}{18.10}}{18.10}}}{Preesquemas étale sobre un preesquema geométricamente unirrama o normal}{157}
\EGAIVcontentsline{{\texorpdfstring{\protect\EGAReadLink{IV.18.11}{18.11}}{18.11}}}{Aplicación a las álgebras locales noetherianas completas sobre un cuerpo}{169}
\EGAIVcontentsline{{\texorpdfstring{\protect\EGAReadLink{IV.18.12}{18.12}}{18.12}}}{Aplicaciones de la localización étale a los morfismos cuasi-finitos (generalizaciones de resultados anteriores)}{181}

\EGAIVcontentsline{\S19}{Inmersiones regulares y planitud normal}{185}
\EGAIVcontentsline{{\texorpdfstring{\protect\EGAReadLink{IV.19.1}{19.1}}{19.1}}}{Propiedades de las inmersiones regulares}{185}
\EGAIVcontentsline{{\texorpdfstring{\protect\EGAReadLink{IV.19.2}{19.2}}{19.2}}}{Inmersiones transversalmente regulares}{190}
\EGAIVcontentsline{{\texorpdfstring{\protect\EGAReadLink{IV.19.3}{19.3}}{19.3}}}{Intersecciones completas relativas (caso plano)}{194}
\EGAIVcontentsline{{\texorpdfstring{\protect\EGAReadLink{IV.19.4}{19.4}}{19.4}}}{Aplicación: criterios de regularidad y de lisitud para los preesquemas obtenidos por explosión}{198}
\EGAIVcontentsline{{\texorpdfstring{\protect\EGAReadLink{IV.19.5}{19.5}}{19.5}}}{Criterios de $M$-regularidad}{204}
\EGAIVcontentsline{{\texorpdfstring{\protect\EGAReadLink{IV.19.6}{19.6}}{19.6}}}{Sucesiones regulares relativamente a un módulo filtrado cociente}{209}
\EGAIVcontentsline{{\texorpdfstring{\protect\EGAReadLink{IV.19.7}{19.7}}{19.7}}}{Criterio de planitud normal de Hironaka}{212}
\EGAIVcontentsline{{\texorpdfstring{\protect\EGAReadLink{IV.19.8}{19.8}}{19.8}}}{Propiedades del paso al límite proyectivo}{219}
\EGAIVcontentsline{{\texorpdfstring{\protect\EGAReadLink{IV.19.9}{19.9}}{19.9}}}{Sucesiones $\sh{F}$-regulares y profundidad}{222}

\EGAIVcontentsline{\S20}{Funciones meromorfas y pseudomorfismos}{223}
\EGAIVcontentsline{{\texorpdfstring{\protect\EGAReadLink{IV.20.0}{20.0}}{20.0}}}{Introducción}{223}
\EGAIVcontentsline{{\texorpdfstring{\protect\EGAReadLink{IV.20.1}{20.1}}{20.1}}}{Funciones meromorfas}{224}
\EGAIVcontentsline{{\texorpdfstring{\protect\EGAReadLink{IV.20.2}{20.2}}{20.2}}}{Pseudomorfismos y pseudofunciones}{231}

\oldpage[IV]{343}
\EGAIVcontentsline{{\texorpdfstring{\protect\EGAReadLink{IV.20.3}{20.3}}{20.3}}}{Composición de los pseudomorfismos}{237}
\EGAIVcontentsline{{\texorpdfstring{\protect\EGAReadLink{IV.20.4}{20.4}}{20.4}}}{Propiedades de los dominios de definición de las funciones racionales}{244}
\EGAIVcontentsline{{\texorpdfstring{\protect\EGAReadLink{IV.20.5}{20.5}}{20.5}}}{Pseudomorfismos relativos}{249}
\EGAIVcontentsline{{\texorpdfstring{\protect\EGAReadLink{IV.20.6}{20.6}}{20.6}}}{Funciones meromorfas relativas}{252}

\EGAIVcontentsline{\S21}{Divisores}{255}
\EGAIVcontentsline{{\texorpdfstring{\protect\EGAReadLink{IV.21.1}{21.1}}{21.1}}}{Divisores sobre un espacio anillado}{255}
\EGAIVcontentsline{{\texorpdfstring{\protect\EGAReadLink{IV.21.2}{21.2}}{21.2}}}{Divisores e ideales fraccionarios invertibles}{258}
\EGAIVcontentsline{{\texorpdfstring{\protect\EGAReadLink{IV.21.3}{21.3}}{21.3}}}{Equivalencia lineal de los divisores}{263}
\EGAIVcontentsline{{\texorpdfstring{\protect\EGAReadLink{IV.21.4}{21.4}}{21.4}}}{Imágenes recíprocas de divisores}{265}
\EGAIVcontentsline{{\texorpdfstring{\protect\EGAReadLink{IV.21.5}{21.5}}{21.5}}}{Imágenes directas de divisores}{267}
\EGAIVcontentsline{{\texorpdfstring{\protect\EGAReadLink{IV.21.6}{21.6}}{21.6}}}{Ciclo 1-codimensional asociado a un divisor}{270}
\EGAIVcontentsline{{\texorpdfstring{\protect\EGAReadLink{IV.21.7}{21.7}}{21.7}}}{Interpretación de los ciclos positivos de codimensión 1 en términos de subpreesquemas}{277}
\EGAIVcontentsline{{\texorpdfstring{\protect\EGAReadLink{IV.21.8}{21.8}}{21.8}}}{Divisores y normalización}{280}
\EGAIVcontentsline{{\texorpdfstring{\protect\EGAReadLink{IV.21.9}{21.9}}{21.9}}}{Divisores sobre los preesquemas de dimensión 1}{284}
\EGAIVcontentsline{{\texorpdfstring{\protect\EGAReadLink{IV.21.10}{21.10}}{21.10}}}{Imágenes recíprocas e imágenes directas de ciclos 1-codimensionales}{289}
\EGAIVcontentsline{{\texorpdfstring{\protect\EGAReadLink{IV.21.11}{21.11}}{21.11}}}{Factorialidad de los anillos regulares}{302}
\EGAIVcontentsline{{\texorpdfstring{\protect\EGAReadLink{IV.21.12}{21.12}}{21.12}}}{El teorema de pureza de Van der Waerden para el conjunto de ramificación de un morfismo birracional}{304}
\EGAIVcontentsline{{\texorpdfstring{\protect\EGAReadLink{IV.21.13}{21.13}}{21.13}}}{Pares parafactoriales. Anillos locales parafactoriales}{313}
\EGAIVcontentsline{{\texorpdfstring{\protect\EGAReadLink{IV.21.14}{21.14}}{21.14}}}{El teorema de Ramanujam--Samuel}{323}
\EGAIVcontentsline{{\texorpdfstring{\protect\EGAReadLink{IV.21.15}{21.15}}{21.15}}}{Divisores relativos}{329}
\EGAIVcontentsline{}{BIBLIOGRAFÍA}{333}
\EGAIVcontentsline{}{ÍNDICE DE NOTACIONES}{334}
\EGAIVcontentsline{}{ÍNDICE TERMINOLÓGICO}{336}

\endgroup
